\documentclass[a4paper]{article}
% Espanol (Mexico) con XeLaTeX: fontspec para fuentes Unicode, polyglossia
% para la division silabica y los nombres del documento en la variante mexicana.
\usepackage{fontspec}
\usepackage{polyglossia}
\setdefaultlanguage[variant=mexican]{spanish}
% polyglossia no traduce los nombres de listings.
\AtBeginDocument{\providecommand{\lstlistingname}{}\renewcommand{\lstlistingname}{Listado}%
  \providecommand{\lstlistlistingname}{}\renewcommand{\lstlistlistingname}{Índice de listados}}
\usepackage{amsmath,amssymb}
\usepackage{graphicx}
\usepackage[margin=2.35cm]{geometry}
\usepackage[most]{tcolorbox}
\usepackage{etoolbox}
\usepackage{listings}
\usepackage{hyperref}
\usepackage{booktabs}
\usepackage{array}
\usepackage{float}
\usepackage{xcolor}

\definecolor{kbBack}{HTML}{F2F6FA}\definecolor{kbFrame}{HTML}{9DB4CC}\definecolor{kbTitle}{HTML}{52708F}
\definecolor{ibBack}{HTML}{FBF5E7}\definecolor{ibFrame}{HTML}{C9A961}\definecolor{ibTitle}{HTML}{7D5F1E}
\definecolor{wbBack}{HTML}{FCF2F0}\definecolor{wbFrame}{HTML}{D6A096}\definecolor{wbTitle}{HTML}{9E4F43}
\definecolor{dbBack}{HTML}{F4F7F3}\definecolor{dbFrame}{HTML}{AEC2AC}\definecolor{dbTitle}{HTML}{5F7F64}

\tcbset{softbox/.style={enhanced,breakable,boxrule=0.8pt,arc=2pt,left=3mm,right=3mm,top=2mm,bottom=2mm,coltitle=white,fonttitle=\bfseries,toptitle=0.6mm,bottomtitle=0.6mm}}
\newtcolorbox{knowledgebox}[1]{softbox,colback=kbBack,colframe=kbFrame,colbacktitle=kbTitle,title={#1}}
\newtcolorbox{importantbox}[1]{softbox,colback=ibBack,colframe=ibFrame,colbacktitle=ibTitle,title={#1}}
\newtcolorbox{warningbox}[1]{softbox,colback=wbBack,colframe=wbFrame,colbacktitle=wbTitle,title={#1}}
\newtcolorbox{dialoguebox}[1]{softbox,colback=dbBack,colframe=dbFrame,colbacktitle=dbTitle,title={#1}}

\lstset{language=Python,basicstyle=\ttfamily\small,keywordstyle=\color{blue!65!black},stringstyle=\color{red!60!black},commentstyle=\color{green!45!black},breaklines=true,frame=single,numbers=left,numberstyle=\tiny\color{gray},captionpos=b,extendedchars=true}
\hypersetup{colorlinks=true,linkcolor=blue!50!black,urlcolor=blue!60!black}
\setlength{\parindent}{2em}
\setlength{\parskip}{0.35em}
\setlength{\emergencystretch}{2em}
\renewcommand{\arraystretch}{1.25}

\newcommand{\makelecturetitle}{
\begin{titlepage}
\centering
\vspace*{0.8cm}
{\huge\bfseries \notetitle\par}
\vspace{0.6cm}
{\large \lecturers\par}
\vspace{0.25cm}
{\large \noteauthors\par}
\vspace{0.8cm}
\includegraphics[width=0.86\textwidth,height=0.44\textheight,keepaspectratio]{cover.jpg}
\vfill
\begin{tcolorbox}[width=0.92\textwidth,colback=black!2!white,colframe=black!55,sharp corners]
\textbf{Curso}: Stanford CS329A Self-Improving AI Agents (Autumn 2025)\par
\textbf{Fecha de la clase}: \lecturedate\quad
\textbf{Fecha de publicación del video}: 2026-08-03\par
\textbf{Fecha de elaboración de las notas}: \notedate\par
\textbf{Canal / duración}: Stanford Online\quad / \videoduration\par
\textbf{Enlace del video}: \href{\videourl}{\nolinkurl{\videourl}}\par
\textbf{Normas de elaboración}: \href{https://github.com/wdkns/wdkns-skills}{wdkns/wdkns-skills} (commit \texttt{39f1a04}) y las normas de calidad de este proyecto
\end{tcolorbox}
\end{titlepage}
\tableofcontents
\newpage
}

\newcommand{\videofigure}[4][0.95]{
\begin{figure}[H]
\centering
\includegraphics[width=#1\textwidth]{#2}
\caption{#3\protect\footnotemark}
\end{figure}
\footnotetext{Intervalo del video: #4.}
}


\newcommand{\notetitle}{Stanford CS329A: Agentes de IA autosuperación\\Part 8: Agentic Evaluation y tareas de largo horizonte}
\newcommand{\lecturers}{Aakanksha Chowdhery}
\newcommand{\noteauthors}{AI Course Notes \& Codex}
\newcommand{\notedate}{10 de agosto de 2026}
\newcommand{\lecturedate}{2025-11-17}
\newcommand{\videoduration}{1:15:17}
\newcommand{\videourl}{https://www.youtube.com/watch?v=8JAqLnTaZu4}

\begin{document}
\makelecturetitle

\section{Por qué la evaluación de agentes de próxima generación no puede basarse solo en la precisión}

Los benchmarks tradicionales comprimen la tarea en una entrada estática y una respuesta corta, adecuados para medir conocimiento y razonamiento local, pero difícilmente reflejan la capacidad del agent de planificar, invocar herramientas, recuperarse de errores y entregar productos complejos a lo largo de decenas de minutos u horas. Esta clase responde tres preguntas mediante tres evaluaciones complementarias:

\begin{enumerate}
    \item \textbf{METR Time Horizon}: ¿cuánto tiempo de una tarea humana puede completar el modelo de forma autónoma con una confiabilidad dada?
    \item \textbf{GDPval}: ¿tiene la salida del modelo, comparada con la de expertos reales de la industria, un valor económico de uso?
    \item \textbf{DeepScholar-Bench}: ante una tarea de investigación abierta, ¿puede el sistema recuperar información de manera exhaustiva, organizar el conocimiento y ofrecer citas verificables?
\end{enumerate}

\videofigure{figures/measuring-ai-progress.jpg}{Evaluar el progreso de la IA requiere considerar simultáneamente la complejidad de la tarea, la confiabilidad y el valor económico; los benchmarks tradicionales se saturan rápidamente.}{00:02:54--00:04:26}

\begin{importantbox}{La capacidad del agent tiene al menos tres ejes independientes}
La duración de la tarea responde «cuánto tiempo puede hacer», la probabilidad de éxito responde «qué tan confiable es» y la calidad de la salida responde «si vale la pena usarlo». Un solo puntaje no puede expresar estos tres ejes a la vez.
\end{importantbox}

\begin{warningbox}{El puntaje de un benchmark no es la proporción real de automatización}
Las tareas experimentales suelen ser más limpias, con un contexto más completo y herramientas más estables, y rara vez incluyen comunicación organizacional, conocimiento tácito y restricciones de responsabilidad. Que un modelo tenga éxito en un benchmark no significa que la misma proporción de puestos reales pueda completarse sin supervisión humana.
\end{warningbox}

\subsection{Resumen de la sección}

Lo que se evalúa en los agents de largo plazo ya no es una respuesta, sino una trayectoria de ejecución y un producto de trabajo final. El tiempo, la confiabilidad, el valor económico, la dependencia del contexto y la verificabilidad deben integrarse juntos en el sistema de métricas.

\section{METR Time Horizon: usar el tiempo de finalización humano para calibrar la dificultad de las tareas}

\subsection{Definición del 50\% time horizon}

\videofigure{figures/time-horizon-definition.jpg}{Time horizon se define como la duración de finalización humana correspondiente a las tareas que el modelo puede completar con una tasa de éxito dada.}{00:04:26--00:06:08}

Denotemos como $t$ el tiempo que necesita un humano competente para completar la tarea, y como $P_\theta(\text{success}\mid t)$ la probabilidad de éxito del modelo. METR usa una logistic curve para ajustar la variación de la tasa de éxito según la duración de la tarea:

$$
P_\theta(\text{success}\mid t)=
\frac{1}{1+\exp\left[\beta\left(\log t-\log H_{50}\right)\right]}.
$$

\begin{itemize}
    \item $t$: el tiempo de finalización de un humano competente;
    \item $H_{50}$: la duración de la tarea correspondiente a que el modelo alcance una tasa de éxito del 50\%;
    \item $\beta$: la pendiente con la que la tasa de éxito disminuye conforme aumenta la duración de la tarea;
    \item $P_\theta$: la probabilidad de éxito empírica bajo una configuración de agent dada.
\end{itemize}

Time horizon no es cuánto tiempo estuvo funcionando el modelo, sino \textbf{cuánto tiempo le tomaría a un humano completar la tarea que el modelo puede completar}. Aquí el tiempo humano funciona como una regla empírica de la complejidad de la tarea.

\begin{importantbox}{El tiempo humano es un proxy de dificultad, no un proxy de valor}
Una tarea mecánica de organización que toma tres horas puede tener un valor muy bajo, mientras que un juicio de alto riesgo que toma diez minutos puede tener un valor muy alto. Time horizon mide la extensión de la ejecución autónoma, y no se puede equiparar directamente con el impacto económico.
\end{importantbox}

\subsection{Conjunto de tareas y proceso de medición}

\videofigure{figures/metr-task-suite.jpg}{METR combina operaciones atómicas del orden de segundos, tareas de software e investigación de minutos a horas, y ML research tasks más largas.}{00:06:08--00:07:08}

\videofigure{figures/metr-method.jpg}{El proceso de evaluación incluye construir tareas de múltiples escalas, registrar la línea base humana, ejecutar el agent, estimar la tasa de éxito y ajustar el time horizon.}{00:07:08--00:11:16}

Las tareas abarcan ingeniería de software, ciberseguridad, razonamiento general e investigación en machine learning. Para cada tarea, primero se deja que un profesional competente la complete de forma independiente y se registra el tiempo, y luego se deja que el agent la intente varias veces en un entorno controlado, con el éxito determinado por pruebas automáticas o reglas manuales.

\begin{knowledgebox}{Por qué se ajusta con log-time}
La duración de las tareas abarca varios órdenes de magnitud: segundos, minutos, horas. Ajustar sobre $\log t$ permite que tareas de distintas escalas se distribuyan de manera más uniforme en la curva, y hace que «duplicar la capacidad» se manifieste como una tendencia aproximadamente lineal.
\end{knowledgebox}

\videofigure{figures/success-vs-human-time.jpg}{La tasa de éxito del modelo tiene una correlación negativa clara con el tiempo de finalización humano, lo que indica que la duración humana puede explicar una parte considerable de la dificultad para el agent.}{00:12:22--00:14:02}

En la figura todavía queda un residual considerable: algunas tareas que toman mucho tiempo pero tienen pasos regulares resultan relativamente fáciles para el agent, mientras que algunas tareas cortas que requieren un juicio fino o la manipulación de herramientas frágiles resultan, por el contrario, difíciles. Por eso el time horizon es una tendencia agregada, no un predictor para una tarea individual.

\subsection{Resumen de la sección}

METR usa el tiempo humano para mapear tareas heterogéneas a un eje de dificultad unificado, y luego lee $H_{50}$ o $H_{80}$ a partir de la curva de tasa de éxito. Es adecuado para dar seguimiento a tendencias de capacidad a largo plazo, pero es necesario conservar el análisis por tipo de tarea y por modo de falla.
\section{El crecimiento del horizonte temporal y la brecha de confiabilidad}

\subsection{De segundos a casi una hora}

\videofigure{figures/fifty-percent-horizon.jpg}{50\% horizon: crecimiento desde el nivel de segundos de GPT-2 hasta casi una hora en Claude 3.7 Sonnet; la tendencia histórica se aproxima a un crecimiento exponencial.}{00:14:02--00:16:18}

El ajuste presentado en el curso muestra que, en los últimos seis años aproximadamente, el doubling time del time horizon ha sido de unos siete meses. Esta tendencia implica que el agent no solo es más preciso en tareas cortas del mismo tipo, sino que también puede mantener cadenas más largas de planificación y llamadas a herramientas.

Si el horizon crece exponencialmente con el año $y$, puede escribirse como:

$$
H(y)=H_0\,2^{(y-y_0)/\tau},
$$

\begin{itemize}
    \item $H_0$: el time horizon en el año de referencia;
    \item $y-y_0$: los años transcurridos;
    \item $\tau$: el doubling time, que el curso estima en unos siete meses.
\end{itemize}

\begin{warningbox}{No tomes la extrapolación de un historial corto como una predicción certera}
La frecuencia de publicación de los modelos, la composición del conjunto de tareas, el entorno de herramientas y la filtración de datos pueden alterar la curva. El ajuste exponencial describe las muestras pasadas; no garantiza que se mantenga la misma pendiente durante los próximos seis años.
\end{warningbox}

\subsection{50\% y 80\% horizon}

\videofigure{figures/eighty-percent-horizon.jpg}{El 80\% success horizon también mejora entre generaciones, pero su duración absoluta es considerablemente más corta que la del 50\% horizon.}{00:19:42--00:21:42}

\videofigure{figures/horizon-comparison.jpg}{El 50\% horizon de Claude 3.7 es de unos 59 minutos, mientras que el 80\% horizon es de unos 15 minutos, lo que revela una brecha de confiabilidad en tareas de dificultad media.}{00:21:42--00:23:18}

La métrica del 50\% responde a «qué tan largo es lo que hay la mitad de probabilidad de lograr»; la métrica del 80\% se acerca más a lo que exige un despliegue real. Las pendientes de crecimiento de ambas pueden ser similares, pero $H_{80}$ es varias veces más corta, lo que indica que el modelo sigue presentando fallas aleatorias inaceptables con frecuencia en tareas más largas.

\begin{importantbox}{Lo que importa para el despliegue es la reliability frontier}
Un sistema real debería reportar una curva $H_p$ completa, no solo $H_{50}$. Las tareas médicas, financieras o de infraestructura pueden requerir $p\ge 0.99$; en cambio, la generación de borradores de bajo riesgo puede aceptar fallas seguidas de un reintento.
\end{importantbox}

\subsection{Resumen de la sección}

El alcance de tareas autónomas del agent crece rápidamente, pero sigue existiendo una brecha enorme entre «lograrlo alguna vez» y «lograrlo de manera constante». Cuanto más alto es el umbral de confiabilidad exigido, más corta es la duración de las tareas que pueden completarse de forma automática.

\section{Modos de fallo de METR y límites de su validez}

\subsection{Por qué fallan las trayectorias largas}

\videofigure{figures/metr-failure-modes.jpg}{Los fallos comunes incluyen una planeación o selección de herramientas inadecuada, razonamiento erróneo, abandono prematuro y ciclos de acciones repetidas.}{00:23:18--00:26:18}

Los errores por lo general no aparecen de golpe en el último paso, sino que se acumulan desde decisiones tempranas: se elige la herramienta equivocada, no se verifican los productos intermedios, se interpreta mal el estado del entorno, y a partir de ahí se sigue invirtiendo tokens y tiempo en la rama equivocada.

\begin{warningbox}{Más context no sustituye la gestión de estado}
Meter todo el historial completo de vuelta al modelo no equivale a que el agent recuerde el estado clave. Las trayectorias largas necesitan un grafo de tareas explícito, hechos ya verificados, problemas sin resolver, resúmenes de las salidas de las herramientas y puntos de retroceso.
\end{warningbox}

\begin{knowledgebox}{El comportamiento cíclico es un fallo de la política de detención}
Que un agent repita la misma acción una y otra vez suele indicar que no registró en su estado que "la acción ya se intentó y no funcionó", ni activó una escalación, un cambio de plan o una terminación. Detectar estados repetidos es más eficaz que simplemente aumentar el número máximo de pasos.
\end{knowledgebox}

\subsection{Si las tareas realmente representan el trabajo real}

\videofigure{figures/metr-validation.jpg}{METR valida sus conclusiones sobre el time-horizon desde perspectivas como el desorden del entorno, el sesgo en la dificultad de las tareas y las revisiones internas.}{00:26:18--00:28:06}

La validación necesita revisar: si la tarea es demasiado limpia, si la calificación automática pasa por alto problemas de calidad, si la estimación del tiempo humano es confiable, si el modelo ya había visto la pregunta, y si los errores del entorno se contabilizaron erróneamente como errores de capacidad. El curso enfatiza que la credibilidad de la evaluación proviene de estas auditorías, no de una bonita línea de tendencia.

\begin{importantbox}{El benchmark también necesita un verifier}
Así como el agent es evaluado por un verifier, el benchmark mismo también debe ser verificado: la representatividad de las tareas, las reglas de calificación, la varianza entre ejecuciones repetidas y las líneas base humanas pueden ser todas fuentes de error sistemático.
\end{importantbox}

\subsection{Resumen de la sección}

METR es adecuado para responder preguntas sobre la amplitud de las capacidades, pero no puede responder por sí solo preguntas sobre la calidad del producto de trabajo, el valor económico y el contexto real. Analizar las trayectorias fallidas y la validez del benchmark es una parte necesaria para interpretar correctamente el time horizon.
\section{GDPval: comparar directamente el output con expertos de la industria}

\subsection{Pasar de "logró completarlo" a "¿la entrega es suficientemente buena?"}

\videofigure{figures/economic-value.jpg}{GDPval no solo pregunta si el agent completó la tarea, sino que hace que expertos comparen la IA con los productos de trabajo reales de profesionales del mismo ramo.}{00:28:06--00:28:38}

\videofigure{figures/gdpval-sectors.jpg}{Las tareas cubren nueve sectores de alto PIB y 44 ocupaciones, en un intento de representar el trabajo de conocimiento común en la economía.}{00:28:38--00:30:02}

\videofigure{figures/gdpval-task-examples.jpg}{Los productos de las tareas incluyen CAD, análisis de inversión, evaluación de imágenes de enfermería, diapositivas, tablas y reportes operativos.}{00:30:02--00:34:16}

A diferencia de los benchmarks de código, los outputs de GDPval por lo general no tienen una única respuesta correcta. La evaluación usa pairwise preference de expertos de la industria: se compara el producto de la IA con el producto del experto humano sin saber cuál es cuál, y se determina si gana, pierde o hay empate.

\begin{importantbox}{GDPval mide la calidad del producto, no el proceso de pensamiento oculto}
Los expertos ven la entrega final, así que la métrica se acerca más a la experiencia del cliente; pero no puede indicar directamente si el proceso del agent es confiable, si citó hechos incorrectos o si fallaría con otra entrada.
\end{importantbox}

\subsection{Construcción de los datos y características de las tareas}

\videofigure{figures/gdpval-construction.jpg}{GDPval primero selecciona los sectores según el PIB, y luego elige las ocupaciones y las actividades de trabajo típicas, con una construcción representativa de arriba hacia abajo.}{00:34:16--00:35:08}

\videofigure{figures/gdpval-task-characteristics.jpg}{El tiempo promedio de finalización humana de una tarea es de aproximadamente siete horas; las tareas son multimodales, dependen de archivos de referencia y, en su mayoría, no tienen una única respuesta correcta.}{00:35:08--00:36:08}

Este tipo de tareas se acerca más al trabajo de oficina que las tareas de software de METR, pero también es más difícil de calificar automáticamente. Una evaluación de alta calidad requiere expertos de dominio, un rubric explícito, pruebas ciegas y verificaciones de consistencia, con un costo mucho mayor que ejecutar pruebas unitarias.

\begin{warningbox}{La evaluación de preferencia se ve afectada por el estilo y la presentación}
Un reporte con buena tipografía y lenguaje fluido puede ocultar errores factuales; un producto humano profesional pero sencillo también puede subestimarse. Se debe registrar por separado accuracy, instruction following, formatting y severity, en lugar de mirar solo la preferencia general.
\end{warningbox}

\subsection{El avance de los modelos no sigue una sola línea}

\videofigure{figures/gdpval-improvement.jpg}{El pairwise win rate en GDPval mejora de manera significativa entre generaciones de modelos, pero de forma escalonada, no como un crecimiento lineal suave.}{00:36:08--00:38:18}

\videofigure{figures/model-specific-strengths.jpg}{Distintos modelos muestran fortalezas claramente diferenciadas: unos destacan en estética y documentos, otros son mejores en cálculo, instruction following o herramientas complejas.}{00:38:18--00:39:08}

El ranking de modelos depende de la combinación de tareas. Un modelo puede dominar en slides, PDF y disposición visual, mientras que otro puede ser más fuerte en cálculo de tablas, precisión textual o herramientas de código. Por eso, el despliegue empresarial debería enrutar por tarea, en lugar de suponer que un solo modelo domina todas las ocupaciones.

\begin{knowledgebox}{El model routing es un producto directo de la evaluación}
Si el benchmark conserva etiquetas de grano fino sobre ocupación, modalidad y tipo de error, se puede entrenar un router: elegir el modelo, la cadena de herramientas y la intensidad de revisión según la tarea, en lugar de usarlo solo para publicar una tabla de posiciones general.
\end{knowledgebox}

\subsection{Severidad de las fallas y costo económico}

\videofigure{figures/gdpval-failure-modes.jpg}{Los problemas comunes incluyen errores de instruction following, formato y precisión; la composición de errores difiere entre modelos.}{00:39:08--00:40:08}

\videofigure{figures/failure-severity.jpg}{Muchas fallas son del tipo "aceptable pero no lo suficientemente bueno"; la proporción de errores verdaderamente catastróficos es baja pero no se puede ignorar.}{00:40:08--00:41:08}

\videofigure{figures/economic-speed-cost.jpg}{Al considerar "intentar varias veces y luego dejar que un humano lo corrija", el modelo puede aportar beneficios de velocidad y costo, pero el beneficio depende del costo de reintento y de revisión.}{00:41:08--00:42:02}

Se puede expresar con un modelo simplificado de costo esperado:

$$
C(n)=nC_{\mathrm{AI}}+
\bigl(1-P_n(\text{acceptable})\bigr)C_{\mathrm{repair}}+C_{\mathrm{review}}.
$$

\begin{itemize}
    \item $n$: número de generaciones o reintentos;
    \item $C_{\mathrm{AI}}$: costo de cada ejecución del modelo;
    \item $P_n$: probabilidad de obtener un resultado aceptable a partir de $n$ intentos;
    \item $C_{\mathrm{repair}}$: costo de que un experto repare un producto fallido;
    \item $C_{\mathrm{review}}$: costo de revisión necesario sin importar si hubo éxito o no.
\end{itemize}

\begin{importantbox}{El beneficio económico proviene del flujo de trabajo "IA + humano"}
Si la revisión y la reparación son baratas, puede haber valor incluso con una tasa de éxito baja en un solo intento; si los errores son difíciles de detectar o tienen consecuencias graves, el costo de la revisión humana anula la ventaja de velocidad del modelo.
\end{importantbox}

\videofigure{figures/performance-variation.jpg}{El desempeño en GDPval varía según el sector, la duración de la tarea y la modalidad; las tareas cortas y algunos sectores se acercan más al nivel experto.}{00:42:02--00:45:48}

\subsection{Resumen de la sección}

GDPval hace avanzar la evaluación de "responder correctamente una pregunta" a "entregar un producto profesional comparable". Revela las fortalezas de los modelos, la severidad de los errores y el costo de la colaboración humano-máquina, pero sigue estando influido por el contexto y las preferencias de los expertos.

\section{Context Problem: el trabajo real no es un prompt completo}

\videofigure{figures/context-experiment.jpg}{Cuando GDPval acorta deliberadamente el prompt y retira parte del contexto, la tasa de victorias del modelo cae claramente, lo que revela la importancia de la context acquisition.}{00:45:48--00:46:48}

\videofigure{figures/context-problem.jpg}{Puede haber una diferencia de eficiencia de varias veces entre quien mantiene el repositorio, un experto del dominio y un contratista externo, porque el trabajo real depende de contexto implícito.}{00:46:48--00:47:32}

Un empleado real conoce las decisiones históricas, las preferencias de la organización, la ubicación de los archivos, los flujos de excepción y quién puede responder qué. Los benchmarks suelen organizar de antemano toda esta información dentro del prompt, con lo cual se salta uno de los trabajos más difíciles del agent: descubrir qué falta, a quién preguntar, dónde buscar y cómo juzgar si la información sigue siendo válida.

\begin{warningbox}{Es fácil confundir la falta de contexto con la falta de capacidad}
Un agent inteligente pero sin acceso al repositorio, a los clientes ni al contexto organizacional puede rendir peor que una persona con capacidad ordinaria pero con contexto suficiente. La evaluación debe distinguir entre los fallos causados por el reasoning, por la context retrieval y por los permisos y el acceso al entorno.
\end{warningbox}

\begin{importantbox}{La próxima generación de agents debe aprender context acquisition}
El agent no debe esperar pasivamente un prompt perfecto, sino enumerar de forma activa la información faltante, recuperar material interno, hacer preguntas de aclaración y mantener el registro de fuentes y vigencia. Esta es la capacidad clave para pasar de ser un benchmark solver a ser un colega real.
\end{importantbox}

\videofigure{figures/gdpval-implications.jpg}{GDPval respalda más bien la asistencia profesional a corto plazo con supervisión humana, y no una sustitución automática e incondicional.}{00:47:32--00:49:10}

\subsection{Resumen de la sección}

El cuello de botella del trabajo real suele estar en la obtención del contexto, no en la generación del texto final. Tanto el despliegue como la evaluación deben incorporar al agent loop la recuperación, los permisos, las aclaraciones y la memoria organizacional.

\section{DeepScholar-Bench: el triple problema de la síntesis de investigación}

\subsection{Por qué el Related Work es una tarea de larga duración ideal}

\videofigure{figures/deepscholar-bench.jpg}{DeepScholar-Bench complementa la duración de tareas de METR y el valor económico de GDPval desde la perspectiva de la síntesis de investigación.}{00:51:48--00:53:04}

\videofigure{figures/deepscholar-task.jpg}{El benchmark exige generar un Related Work a partir del tema de un artículo, lo cual requiere descubrir, filtrar, organizar y citar la investigación relevante.}{00:53:04--00:54:52}

El Related Work no tiene una única respuesta correcta, pero sí una estructura de calidad clara: debe cubrir los trabajos clave, organizarse por temas, describir con precisión las contribuciones, evitar omisiones o invenciones, y permitir que las citas se puedan verificar. Esto lo convierte en un escenario de evaluación de alto valor para los agentes de investigación profunda.

\begin{knowledgebox}{Research synthesis no es una generación de texto largo común}
Incluye recall en la recuperación, juicio sobre las fuentes, alineación claim-to-source, síntesis entre documentos y organización narrativa. Si falla cualquiera de estos eslabones, el párrafo final puede «leerse bien, pero en realidad ser incompleto o no verificable».
\end{knowledgebox}

\subsection{Evaluación tridimensional}

\videofigure{figures/three-dimensional-eval.jpg}{La evaluación cubre simultáneamente la síntesis de conocimiento, la calidad de la recuperación y la verificabilidad de las citas.}{00:54:52--00:56:18}

Las tres dimensiones de la métrica responden respectivamente:

\begin{itemize}
    \item \textbf{Knowledge synthesis}: si la estructura es coherente, si los hechos son correctos, si los conceptos clave se organizan en comparaciones con sentido;
    \item \textbf{Retrieval quality}: si las fuentes son relevantes, exhaustivas e importantes, en lugar de solo encontrar algunos artículos fáciles de buscar;
    \item \textbf{Verifiability}: si las citas existen, si respaldan el claim correspondiente, si el origen se puede rastrear.
\end{itemize}

\videofigure{figures/deepscholar-results.jpg}{Los distintos sistemas destacan cada uno en métricas individuales, pero ninguno supera aproximadamente el 19\% en todas las métricas.}{00:56:18--00:57:26}

\begin{importantbox}{El eslabón más débil determina la credibilidad del informe de investigación}
Una organización buena pero con fuentes incompletas produce un sesgo con buena apariencia; muchas fuentes pero con claims que no corresponden a las citas produce una acumulación no verificable. Es necesario optimizar conjuntamente las tres dimensiones de la métrica.
\end{importantbox}

\subsection{Análisis de fallas}

\videofigure{figures/deepscholar-failures.jpg}{Las principales fallas provienen de la cobertura insuficiente de fuentes, el juicio ineficiente de importancia, la calidad inestable de la síntesis y la falta de hechos citados.}{00:57:26--01:02:10}

Los modelos suelen ser capaces de encontrar artículos «relevantes», pero les cuesta encontrar artículos \textbf{exhaustivos y clave}; también pueden gastar una gran cantidad de tokens revisando repetidamente fuentes marginales. Incluso cuando se recupera la literatura correcta, el sistema todavía debe identificar la contribución, las limitaciones y las relaciones de cada artículo, y por último vincular con precisión el claim a la evidencia.

\begin{warningbox}{Una citation precision alta no implica una coverage alta}
Que cada cita sea real y respalde el claim solo indica que no hay invenciones evidentes; si se omiten los trabajos clave de medio campo, la síntesis sigue sin ser aceptable. Es necesario reportar simultáneamente precision, recall y la cobertura ponderada por importancia.
\end{warningbox}

\begin{knowledgebox}{La condición de detención de la búsqueda es un problema abierto}
El agente de investigación no sabe «qué otros artículos importantes no se han encontrado». Se puede usar el mapa de cobertura de temas, la ganancia de la red de citas, la tasa de repetición de fuentes y la ganancia de hechos nuevos para decidir conjuntamente cuándo detenerse, en lugar de fijar diez búsquedas de antemano.
\end{knowledgebox}

\subsection{Resumen de la sección}

DeepScholar-Bench muestra que el principal cuello de botella de los agentes de investigación de larga duración no es escribir un párrafo fluido, sino la recuperación con alto recall, el juicio de importancia, la síntesis entre documentos y la claim-level verification.
\section{Armar un panorama completo de capacidades a partir de tres tipos de evaluación}

\videofigure{figures/long-horizon-implications.jpg}{METR, GDPval y DeepScholar-Bench exponen, respectivamente, distintos cuellos de botella: el horizonte temporal, el producto económico y la síntesis de investigación.}{01:02:10--01:04:08}

\videofigure{figures/benchmark-realwork-gap.jpg}{Los tres tipos de benchmark presentan diferencias sistemáticas frente al trabajo real: calificación automática, tareas de una sola vez, ausencia de contexto o restricciones de fuentes.}{01:04:08--01:05:42}

\begin{table}[H]
\centering
\small
\begin{tabular}{>{\raggedright\arraybackslash}p{2.6cm}>{\raggedright\arraybackslash}p{3.4cm}>{\raggedright\arraybackslash}p{3.8cm}>{\raggedright\arraybackslash}p{3.5cm}}
\toprule
\textbf{Evaluación} & \textbf{Pregunta central} & \textbf{Principal ventaja} & \textbf{Principal punto ciego} \\
\midrule
METR & ¿Qué tan largas son las tareas que puede completar de forma estable? & Calibración uniforme en tiempo humano, permite rastrear tendencias & Sesgado hacia tareas de software, débil en calidad y valor \\
GDPval & ¿El producto supera al de un experto? & Trabajo real multimodal, interpretación económica directa & La preferencia de expertos es costosa, el contexto se simplifica \\
DeepScholar & ¿Puede realizar una síntesis de investigación confiable? & Evalúa a la vez recuperación, síntesis y citación & Dominios limitados, la exhaustividad sigue siendo difícil de juzgar automáticamente \\
\bottomrule
\end{tabular}
\end{table}

\videofigure{figures/known-unknown.jpg}{Tenemos más certeza de que el modelo progresa rápido en tareas cortas, claras y verificables, pero seguimos sin confianza en entornos de mucho contexto, entornos abiertos y recuperación exhaustiva.}{01:05:42--01:07:18}

\videofigure{figures/evaluation-takeaways.jpg}{El resumen del curso enfatiza la confiabilidad, la recuperación de errores, las características de la tarea y las métricas complementarias, en lugar de perseguir una sola tabla de posiciones.}{01:07:18--01:09:12}

\begin{importantbox}{La evaluación debe convertirse en la interfaz de entrenamiento del agent}
Las etiquetas de fallo de grano fino pueden guiar directamente la mejora: los errores de planeación alimentan los datos del planner, los errores de herramientas alimentan la tool policy, las fuentes omitidas alimentan el retrieval curriculum y las inconsistencias de citación alimentan el entrenamiento del verifier. Un buen benchmark no solo produce un ranking, también debe localizar la señal para la siguiente ronda de auto-mejora.
\end{importantbox}

\subsection{Preguntas abiertas}

\begin{itemize}
    \item ¿Cómo medir la colaboración organizacional a lo largo de días y entre personas, en lugar de una sola tarea en sandbox?
    \item ¿Cómo incorporar el contexto implícito, los permisos, la comunicación y las restricciones de cumplimiento en una evaluación reproducible?
    \item ¿Cómo evaluar la recuperación del agent tras un fallo, y no solo el éxito o fracaso final?
    \item ¿Cómo detectar el reward hacking, la contaminación de benchmarks (benchmark contamination) y el sesgo del evaluador (evaluator bias)?
    \item ¿Cómo combinar la duración de la tarea, la calidad, el costo, el consumo de energía y el riesgo en una frontera de Pareto que permita tomar decisiones?
\end{itemize}

\subsection{Resumen de la sección}

Ningún benchmark único puede representar a un agent de propósito general. Un panorama completo requiere dimensiones complementarias: el horizonte de la tarea, la confiabilidad, la calidad del producto, el costo económico, la obtención de contexto, la cobertura de fuentes y la recuperación de errores.

\section{Lecturas adicionales}

\begin{itemize}
    \item \href{https://metr.org/}{METR: AI task time horizon y estudios de evaluación de agentes de horizonte temporal largo}
    \item \href{https://openai.com/index/gdpval/}{GDPval: evaluación de modelos en tareas económicas reales}
    \item \href{https://arxiv.org/search/?query=DeepScholar-Bench\&searchtype=all}{\emph{DeepScholar-Bench}: Generative Research Synthesis Evaluation}
\end{itemize}

\section{Síntesis y ampliación}

Esta sesión presenta un marco de tres niveles que va de la accuracy en benchmarks a la agent evaluation:

\begin{enumerate}
    \item \textbf{Capability span}: METR usa el tiempo estimado de un humano para completar la tarea con el fin de medir cuánto tiempo puede sostener un modelo una tarea dada una tasa de éxito fija;
    \item \textbf{Economic usefulness}: GDPval hace que expertos de la industria comparen directamente los productos de trabajo de la IA con los profesionales, considerando además el costo de reintentos, correcciones y revisión;
    \item \textbf{Epistemic quality}: DeepScholar-Bench examina la cobertura de fuentes, la síntesis de conocimiento y la citation verifiability.
\end{enumerate}

La conclusión más importante no es que "el agent ya puede hacer tareas de una hora" ni que "se acerca a la tasa de victoria de un experto", sino que el crecimiento de la capacidad es fuertemente condicional:

\begin{itemize}
    \item una tarea larga con una tasa de éxito del 50\% ve su duración reducirse de manera notable al exigir una confiabilidad del 80\%;
    \item el desempeño en un benchmark con un prompt completo no representa la context acquisition en un entorno real;
    \item la calidad visual, el seguimiento de instrucciones, la exactitud factual y la cobertura de citas pueden ser dominados por modelos distintos entre sí;
    \item una salida "aceptable pero no perfecta" puede crear valor en escenarios donde el costo de corrección es bajo, pero no es adecuada para despliegues de alto riesgo sin supervisión;
    \item el progreso en horizontes largos debe ir acompañado de una mejor gestión de estado, políticas de detención, recuperación de errores y verifiers.
\end{itemize}

En definitiva, la evaluación no es una tabla de posiciones única, sino un sistema de diagnóstico para el entrenamiento, el enrutamiento y el control de riesgos. La siguiente sesión regresará al tema central del curso: cómo superar los cuellos de botella de datos, verificación y eficiencia en la automejora, explorando el multi-agent fine-tuning, la meta-verification, Absolute Zero y la intelligence per watt.

\end{document}
